% Section fragment; input from the main document.
\section{模型假设}
\label{sec:frontmatter-assumptions}

为突出购电与储能协同决策的主要矛盾，在题设条件外作如下假设：

\begin{enumerate}
  \item 假设以10分钟区间平均功率表示供需，购电及充放电量可连续调节，忽略设备启停时间与最小动作幅度。
  \item 假设储能充放电效率保持不变，除充放电损耗外，忽略线路损耗、自放电、寿命衰减及运维费用。
  \item 假设外网购电不另设功率上限，微网不向外网售电；无法利用的剩余电量允许弃置，无额外处置费用。
  \item 假设同季节、运行特征相近的历史日可近似反映待决策日的条件分布；情景构造仅使用决策时已知信息。
  \item 假设紧急购电仅补足当前负载缺口，不用于储能充电；问题三、四可在当前区间内依据已观测供需及时调节。
  \item 问题三及其波动电价扩展中，假设按最终生效购电量相对0:00原计划结算一次调整费用，已执行时段不追溯修改；逐次结算另作敏感性分析。
\end{enumerate}
